\section{Cómo codifica la evolución los deseos de alto nivel}

\subsection{Un misterio sin resolver}

Ilya planteó una pregunta que él considera muy misteriosa: ¿cómo codifica la evolución los \textbf{deseos sociales de alto nivel} en el genoma?

\begin{knowledgebox}{Deseos de bajo nivel frente a deseos de alto nivel}
\begin{itemize}
    \item \textbf{Bajo nivel}: el olor de la comida $\rightarrow$ hambre. Es fácil imaginar cómo lo logra la evolución: las señales químicas se conectan directamente con las neuronas de dopamina
    \item \textbf{Alto nivel}: nos importa el prestigio social, lo que piensan los demás, el estatus dentro del grupo. Esto requiere que el cerebro realice una integración de información compleja para entender la «situación social»; ¿cómo especifica la evolución «que te importe el resultado de este cálculo complejo»?
\end{itemize}
\end{knowledgebox}

Ilya propuso una conjetura: quizás la evolución aprovecha la ubicación fija de las regiones cerebrales, las «GPS coordinates of the brain», para especificar «cuando se activan las neuronas de esta ubicación, esto es lo que debe importarte». Pero él mismo señaló un contraejemplo: en los niños a quienes se les extirpa un hemisferio, todas las regiones cerebrales se vuelven a mapear en el hemisferio restante, pero los deseos sociales parecen no verse afectados. Así que esta teoría podría estar equivocada.

\begin{importantbox}{Implicaciones para el alineamiento}
Si la evolución logra, de manera confiable, que los humanos se preocupen por el estatus social y la identidad de grupo incluso en un mundo moderno completamente distinto del entorno original, esto en sí mismo es un \textbf{caso de alineamiento exitoso}. Entender cómo lo logra la evolución podría ofrecer una perspectiva clave para el alineamiento de la IA.
\end{importantbox}

\subsection{Resumen de la sección}

Cómo codifica la evolución los deseos sociales de alto nivel en el genoma es un misterio profundo sin resolver. Entender este mecanismo podría ser crucial para el alineamiento de la IA.


